\documentclass[10pt,a4paper]{amsart}
\usepackage[margin=19mm]{geometry}
\usepackage{amsmath,amssymb,mathtools}
\usepackage[scr=rsfs,cal=euler]{mathalfa}
\usepackage{xfrac}
\usepackage[unicode,hidelinks,bookmarks=false]{hyperref}
\newcommand{\ie}{i.e.\ }
\newcommand{\cf}{cf.\ }
\newcommand{\resp}{respectively\ }
\newcommand{\Subsection}{\subsection}
\providecommand{\nobreakdash}{\nobreak\hbox{-}}
\setlength{\emergencystretch}{2em}
\multlinegap0pt
\usepackage{bm}
\usepackage{bbm}
\usepackage{dsfont}
\usepackage{xspace}
\usepackage{cancel}
\usepackage{mathtools}
\usepackage{amsthm}
\makeatletter
\@ifundefined{lemm}{\newtheorem{lemm}{Lemma}}{}
\@ifundefined{prop}{\newtheorem{prop}[lemm]{Proposition}}{}
\makeatother
\def\Z{{\bf Z}}
\def\lll{\ensuremath{\mathsf l}}
\def\mm{\ensuremath{\mathsf m}}
\title[Numerical invariants of hyper-K\"ahler manifolds]{Numerical invariants of hyper-K\"ahler manifolds}
\author{Olivier Debarre, Chen Jiang}
\date{Source: arXiv:2506.04177v1 (CC BY 4.0)}
\begin{document}
\maketitle

\noindent\emph{Source and licence.} Numerical invariants of hyper-K\"ahler manifolds. Version arXiv:2506.04177v1 (CC BY 4.0), distributed under the Creative Commons Attribution 4.0 International licence, as recorded in the arXiv licence metadata for this exact version and shown on the abstract page https://arxiv.org/abs/2506.04177v1. Full rights notice: licenses/arxiv-2506.04177-CC-BY-4.0.txt.\par\smallskip

\noindent\textit{Editorial scope.} Counted unit: the Prop whose source label is \texttt{prop45} in the article (source0.tex lines 1059--1074, source label \texttt{prop45}), delivered together with the article's own complete original proof of it (source0.tex lines 1076--1133, one \texttt{proof} environment of 58 lines). No mathematics is written, repaired or re-derived: statement and proof are the article's own text line for line. Required context retained uncounted: prr3 (lines 713--719); cxa (lines 901--903); mxroot (lines 908--910); le1 (lines 913--920); le2 (lines 932--940). Every definition, display and label of those lines is retained unchanged. Omitted from this package and not redistributed: 1--712, 720--900, 904--907, 911--912, 921--931, 941--1058, 1075--1075, 1134--1458. Nothing inside a retained excerpt is omitted or altered: every retained body is delivered exactly as the article's lines are written, so no omission receipt of any kind accompanies this package. Citations: the retained excerpts carry no citation command at all, so no bibliography entry of the article is transcribed and no reference is redistributed. Reference adaptation: none was needed and none was made; every reference inside the delivered unit resolves inside it, so no unresolved reference or citation is produced. Printed slips kept literal: no printed word, symbol, hypothesis or number of the counted statement or of its proof was altered. Layout and class: the article's own class is \texttt{amsart} with packages including amsfonts, amssymb, enumerate, bm, hhline, makecell, array, mathtools, mathrsfs, comment, ulem, hyperref, tikz-cd, amscd; the wrapper is the lane's pinned \texttt{amsart} editorial wrapper at 10pt on A4, which restates the theorem-like environments the retained text needs and adds the article's own macro definitions that the retained text needs (\texttt{\textbackslash Z}, \texttt{\textbackslash lll}, \texttt{\textbackslash mm}), with extra wrapper packages (bm, bbm, dsfont, xspace, cancel, mathtools). The delivered numbering is the unit's own. Assets: no retained span contains a figure environment, a graphics-inclusion command, a PDF-page inclusion, a table or a TikZ picture; no image file is copied into this package, the package declares no asset dependency and no image file is redistributed with it. Licence: version arXiv:2506.04177v1 of this article is distributed under the Creative Commons Attribution 4.0 International licence, as recorded in the arXiv licence metadata for this exact version and shown on the abstract page https://arxiv.org/abs/2506.04177v1. A full rights notice is delivered at licenses/arxiv-2506.04177-CC-BY-4.0.txt. Characters: every retained excerpt is pure ASCII, so the delivered engine, XeLaTeX with its default Latin Modern text font, reports no missing character.
\begin{equation}\label{prr3} 
%P_{RR,X}(T)=  \frac{c_X}{720}T^3+\frac{c_Xn_X}{240}T^2+\Bigl(\frac{4}{n_X}+\frac{c_Xn_X^2}{360}\Bigr)T+4.
%\begin{aligned}
P_{RR,X}(T)= \frac{c_X}{720}(T+n_X)^3 +\Bigl(\frac{4}{n_X}-\frac{c_X}{720}n_X^2\Bigr) (T+n_X).
%Q_{RR,X}(T)&= \frac{c_Xn_X^3}{5760}(T+2)^3 +\Bigl(2-\frac{c_Xn_X^3}{1440}\Bigr) (T+2).
%\end{aligned}
\end{equation}

\begin{equation} \label{cxa}
c_Xq_X(\lll,\mm)^n=a(\mm)\frac{(2n)!}{2^nn!}=a(\mm)(2n-1)!!.
\end{equation}

\begin{equation}\label{mxroot}
m_X < 2q_X(\lll,\mm) \sqrt[\leftroot{-3}\uproot{16}  {\scriptstyle {n}}]{\frac{  n!}{a(\mm) }}.
\end{equation}

\begin{lemm}\label{le1}
We have
$$
n! q_X(\lll,\mm)^n \mid a(\mm)C_n$$
and,  if $q_X$ is not even,
$$n! 2^n q_X(\lll,\mm)^n \mid a(\mm)C_n.$$

\end{lemm}

\begin{lemm}\label{le2}
 We have 
\begin{equation*}\label{eq}
a(\mm) \Bigl( \frac{q_X(\mm)+n_X}{2q_X(\lll,\mm)  }-\frac{n-1}2   \Bigr)     \in \Z.
\end{equation*}
In particular,
$$n_X\in   \Z+ \frac{2q_X(\lll,\mm) }{a(\mm)}\Z $$
so that $n_X $ is an integer when $ a(\mm)\in \{1,2\}$.
\end{lemm}

\begin{prop}\label{prop45}
Assume $n=3$ and $a(\mm)=2$.\ Then,
%\begin{enumerate}[{\rm (a)}]
%\item If $n=2$, we have $q_X(\lll,\mm)=1$, $c_X=6$, 
% $n_X\in \{1,2,3\}$, 
%and the quadratic form $q_X$ is even.\ One also has $P_{RR,X}(T)=\frac14T^2+\frac{n_X}{2}T+3$.
%\item If $n=3$, 
 $q_X(\lll,\mm)=1$, $c_X=30$, $n_X \in\{1,2,3,4 \}$, and the quadratic form $q_X$ is even.\ One also has $P_{RR,X}(T)=
   \frac{1}{24}T^3+\frac{n_X}{8}T^2+
 \bigl(\frac{4}{n_X}+\frac{n_X^2}{12}\bigr)T
 % \frac{13}6T
  +4$
%\end{enumerate}
%In both cases, 
and the sublattice $\Z\lll\oplus \Z\mm$ of $(H^2(X,\Z),q_X)$ is a hyperbolic plane.
\end{prop}

\begin{proof}
% \noindent (a) Assume $n=2$.\ We have $C_2=12$ and we obtain   from Lemma~\ref{le1}
%\begin{equation*} 
%q_X(\lll,\mm)^2\mid 12  \qquad\textnormal{(and $ q_X(\lll,\mm)^2 \mid 3 $ if $q_X$ is not even)},
%\end{equation*}
%so that $ q_X(\lll,\mm)\in\{1,2\}$ (and $ q_X(\lll,\mm)=1 $ if $q_X$ is not even).
%
%\medskip
%
%\noindent{\bf{Assume  $ q_X(\lll,\mm)=1 $.}}\ We have $c_X=6$ from~\eqref{cxa}, Lemma~\ref{le2} gives $n_X\in  \Z$, and~\eqref{mxroot} gives $m_X<2 $, so that $n_X \in\{1,2,3\}$.\ Furthermore, we have, by~\eqref{prr2},
%$$P_{RR,X}(T)=\frac14T^2+\frac12n_XT+3, 
%$$
%so that $\frac14 q^2+\frac12n_Xq=\frac14 q(q+2n_X)$ must be an integer for all values $ q$ taken by $q_X$.\ This implies that $q$ must be even.
%
%\medskip
%
%\noindent{\bf{Assume  $ q_X(\lll,\mm)=2 $.}}\ The quadratic form $q_X$ is even, we have $c_X=\frac32$ from~\eqref{cxa}, Lemma~\ref{le2} gives $m_X\in  \Z$, and~\eqref{mxroot} gives $m_X<4 $, so that $m_X \in\{1,2,3\}$.\ As above, $\frac1{16} (2q)^2+\frac18n_X(2q)=\frac14 q(q+2m_X)$ must be an integer for all values $ 2q$ taken by $q_X$.\ This implies that $q$ must be even and contradicts the fact that the gcd of all integers $2q$ represented by $q_X$ is $2$.\ So this case does not happen. 
%
%\bigskip
% \noindent (b) Assume $n=3$.\ 
 We have $C_3=2^5\cdot 3^3\cdot 5$ and we obtain   from Lemma~\ref{le1}
\begin{equation*} 
q_X(\lll,\mm)^3\mid    2^5\cdot 3^2\cdot 5 \qquad\textnormal{(and $ q_X(\lll,\mm)^3 \mid 2^2 \cdot 3^2\cdot 5 $ if $q_X$ is not even),}
\end{equation*}
so that $ q_X(\lll,\mm)\in\{1,2\}$ (and $ q_X(\lll,\mm)=1 $ if $q_X$ is not even).

\medskip

\noindent{\bf{Assume  $ q_X(\lll,\mm)=1 $.}}\ We have $c_X=30$ from~\eqref{cxa}, Lemma~\ref{le2} gives $ n_X\in \Z$, and~\eqref{mxroot} gives $m_X<2\sqrt[3]{3}\sim 2.9 $, so that $n_X \in\{1,2,3,4,5\}$.\ Furthermore, we have, by~\eqref{prr3}, 
$$P_{RR,X}(T)=\frac{1}{24}(T+n_X)^3+\Bigl(\frac{4}{n_X}-\frac{1}{24}n_X^2\Bigr)(T+n_X),
$$
 For all values $q$ taken by $q_X$, this must be an integer when $T=q$, so that
\begin{equation*}%\label{eqdiv3c}
24n_X\mid   n_X(q+n_X)^3+ (96- n_X^3 )(q+n_X).
\end{equation*}
 In particular, $n_X\mid 96 q$.\ If $n_X=5$, this implies $n_X\mid  q$, which is impossible because   the gcd of all integers $q$ represented by $q_X$ is either $1$ or $2$.\ Otherwise, $8n_X\mid 96$, hence we obtain
\begin{equation}\label{eqdiv3cc}
8 \mid  (q+n_X)^3 - n_X^2  (q+n_X)=q(q+n_X)(q+2n_X).
\end{equation}
\begin{itemize}
\item When $n_X=1$, the relation~\eqref{eqdiv3cc} is equivalent to  $q\equiv 0,2,4,6,7\pmod{8}$; this means that every odd value taken by $q_X$ is $\equiv  7\pmod{8}$.\ Assume there exists $\alpha$ such that    $q_X(\alpha)\equiv  7\pmod{8}$.\ Since $q_X(k\lll+\mm)= 2k+q_X(\mm)$, the integer $q_X(\mm)$ must be even and we may even assume $q_X(\mm)=0$.\    For all   $t,u\in\Z$, the integer
$$q_X(t\lll+u\mm+\alpha)=2tu+2tq_X(\lll,\alpha)+2uq_X(\mm,\alpha)+q_X(\alpha)
$$
is odd, hence is $\equiv  7\pmod{8}$.\ This implies $tu+ tq_X(\lll,\alpha)+ uq_X(\mm,\alpha)\equiv 0\pmod4$.\ Taking $t=1$ and $u=0$, we obtain $q_X(\lll,\alpha) \equiv 0\pmod4$; taking $t=0$ and $u=1$, we obtain $q_X(\mm,\alpha) \equiv 0\pmod4$; taking $t=u=1$, we obtain a contradiction.\ Hence $q\equiv 0,2,4,6\pmod{8}$ and $q_X$ is even.\ 
\item When $n_X=2$, the relation~\eqref{eqdiv3cc} is equivalent to  $q\equiv 0,2,4,6\pmod{8}$ and $q_X$ is even.
\item When $n_X=3$, the relation~\eqref{eqdiv3cc} is equivalent to  $q\equiv 0,2,4,5, 6\pmod{8}$.\ If the case $q\equiv 5\pmod{8}$ occurs, the same reasoning as in the case $n_X=1$,  $q\equiv 7\pmod{8}$, gives   a contradiction, hence $q_X$ is even.
\item When $n_X=4$, the relation~\eqref{eqdiv3cc} is equivalent to  $q\equiv 0,2,4,6\pmod{8}$ and $q_X$ is even.
 \end{itemize}

\medskip

\noindent{\bf{Assume  $ q_X(\lll,\mm)=2 $.}}\ The quadratic form $q_X$ is even, we have $c_X=\frac{15}4$ from~\eqref{cxa}, Lemma~\ref{le2} gives $m_X\in   \Z$, and~\eqref{mxroot} gives $m_X<5.8 $, so that $m_X \in\{1,2,3,4,5\}$.\  As above, we  deduce from~\eqref{prr3} that
 $$ \frac{1}{8\cdot 24}(2q+ 2m_X)^3+\Bigl(\frac{2}{m_X}-\frac{1}{8\cdot 24}4m_X^2\Bigr)(2q+ 2m_X) $$
 must be an integer for all values $2q$ taken by $q_X$, so that 
$$24 m_X\mid    m_X(q+m_X)^3+ (96- m_X^3 )(q+m_X).
$$
We reason as above to conclude that the integer~$q$ must be even,  so that all values taken by the quadratic form~$q_X$ are divisible by $4$.\ This is   impossible because   the gcd of all values taken  by $q_X$ is  $2$.\ So this case does not occur. 
 \end{proof}

\end{document}
